\begin{resumo}
A ocorrência de eventos adversos (EAs) no contexto hospitalar representa um dos mais críticos desafios globais de saúde pública, impactando diretamente a morbimortalidade de pacientes e gerando expressivo ônus socioeconômico aos sistemas de saúde. Tradicionalmente, as instituições dependem de sistemas de notificação voluntária, que capturam menos de 10\% dos danos reais. Em resposta a essa lacuna, o \textit{Institute for Healthcare Improvement} (IHI) desenvolveu a metodologia \textit{Global Trigger Tool} (GTT), baseada na revisão ativa e retrospectiva de prontuários por meio de rastreadores ou ``gatilhos'' clínicos. O presente trabalho apresenta a proposta do \textbf{SIGEA-GTT} (\textit{Sistema Inteligente de Gestão de Eventos Adversos – Global Trigger Tool}), uma plataforma computacional interdisciplinar concebida em cooperação entre o Departamento de Computação (DCOMP) e o Departamento de Enfermagem da Universidade Federal de Sergipe (UFS). O sistema integra a metodologia IHI-GTT (6 módulos especializados, 53 gatilhos clínicos e categorização de dano NCC MERP de A a I) a ferramentas consagradas de engenharia e gestão da qualidade hospitalar (Diagrama de Ishikawa, Matriz GUT, Plano de Ação 5W2H e Ciclo PDCA/PDSA). A arquitetura proposta é orientada a microserviços/módulos independentes com backend em Java 21 LTS e Spring Boot 3.3, persistência em PostgreSQL 16 com suporte a JSONB, e frontend responsivo em Angular 18 com Standalone Components, Signals e TailwindCSS, otimizado para estações de trabalho clínicas com alta densidade informacional. O projeto contempla ainda um módulo educacional inovador com prontuários simulados para capacitação profissional e acadêmica no Hospital Universitário da UFS (HU-UFS), submetido a um rigoroso critério de qualidade de software (\textit{Quality Gate}) com 100\% de cobertura de testes.

\textbf{Palavras-chave}: Segurança do Paciente. Global Trigger Tool. Eventos Adversos. Gestão da Qualidade. Engenharia de Software Hospitalar. Universidade Federal de Sergipe.
\end{resumo}
